\paragraph{The integral comparison away from the distinguished prime.}
\label{par:off-prime-row-trace}
\hypertarget{source-paired-off-prime-trace}{}
Fix the run used at \(v\) and keep its base weight, base
character and half level. Denote its distinguished rational
prime by \(p_\circ\). In the argument let \(\lambda\) be any
finite place of \(K\) of residue characteristic different
from \(p_\circ\). The character of the finite trace is the
one on the arithmetic row, namely \(\psi_0^{-1}\) in the
base tuple convention. A generator of the source trace
lattice is
\[
 t=\sum_{u\in U}\psi(u)\operatorname{cusp}(f_u).
\]
Here \((f_u)_{u\in U}\) is one orbit family of short theta
rules.

The orbit family has a common unimodular source lattice, an
admissible odd level denominator, a common harmonic
polynomial and a common residue datum. The labels of the
rules are obtained from that datum by the prescribed level action. The target module \(J\) is the integral span, in
these cusp coordinates, of the base kernel traces. An
admissible kernel trace uses a finite split-supported
source translate and a coefficient in its integral
harmonic lattice. Its form coordinate is specified by
evaluation at the lowest target types. We prove that
\(t\in J_\lambda\).
\label{cor:paired-off-prime-trace-1}

\begin{proposition}[Integral kernel trace away from the distinguished prime]
\label{prop:off-prime-trace-inclusion}
At each place of residue characteristic different from
\(p_\circ\), the orbit trace lattice maps into the kernel theta
lattice in the common form coordinates. For the fixed coefficient character on all finite vectors, \eqref{eq:off-prime-trace-identity} reads
\[
 \mathcal A(t')(\mathbf w)=\sigma(c)\,T_\psi^{\rm ad}S_{\Phi_{\rm fr}}(h,y;\mathbf w)\quad(\text{all }\mathbf w),\qquad v(c)=0.
\]
 The inclusion also holds after the good
projector, on the same local row.
\end{proposition}

\begin{proof}
\proofstep{Integral trivialization on the adapted orbit.}
First adapt the lattice of the orbit to \(\lambda\). The
congruence adjustment of a self-dual lattice makes it
standard locally over the residue characteristic of
\(\lambda\). It preserves the standard lattice at the
old bad and level primes and is the identity modulo the
level there. In the second adjustment, write its finite
adelic lattice translate as \(\gamma h k\). Here \(\gamma\)
is rational unitary, \(h\) has split support away from
the forbidden primes and is identity at infinity, and
\(k\) preserves the standard lattice and its congruence
data. For this factorization use rational unitary
approximation with a common denominator supported
only at split primes outside the prescribed set.
Write \(g\) for the adjusted finite matrix after removal
of \(\gamma\).

Let \(e\) be the finite-adelic idempotent of
its allowed split support, including both conjugate halves.
Set
\[
 h=eg+(1-e)1,\qquad k=e1+(1-e)g.
\]
Since \(e(1-e)=0\) and \(\bar e=e\), the products split into
their two idempotent components. Writing \(*\) for the unitary involution,
\[
 hk=eg+(1-e)g=g,\qquad
 h^*h=e\,g^*g+1-e=1,\qquad k^*k=e+(1-e)g^*g=1.
\] Off the allowed support \(g\) lies in the integral
congruence stabilizer, so \(k\) preserves its lattice and
level marking.
This factorization is used for the conjugated
lattice because the coefficient convention contracts
through conjugate source coordinates. Thus the adjusted
polynomial lattice is the conjugate of the lattice read
at \(h\). The representative has a common denominator
prime to the old bad and level primes and to the residue
characteristic of \(\lambda\). The two rational unitary
changes push forward the polynomial and keep the residue
label at the level.

The invariance of the harmonic
equations and the short theta transformation with even
source rank show that all the short rules change by the
same root of unity \(\zeta\in K_0\).
The common scalar leaves the coefficient valuation unchanged, since
\[
  \zeta^n=1\ \text{for some }n\ge1\quad\Longrightarrow\quad v_\lambda(\zeta)=0.
\]
 The common factor
comes from the normalized theta multiplier, which is
independent of the translated residue label. The multiplier \(\zeta\) is in the base field \(K_0\), so it factors out of the trace on the whole orbit. For the adapted orbit and its trace one obtains
\[
        f'_u=\zeta f_u,\qquad t'=\zeta t,\qquad
        \mathcal O_{K,\lambda}\,t'=\mathcal O_{K,\lambda}\,t .
\]
The last equality uses that both \(\zeta\) and its inverse are integral. The adjustment can
increase the level denominator away from the indicated
primes, where it has no effect on the local inclusion.
\label{cor:paired-off-prime-trace-2}

For the adapted polynomial \(P\), choose a rational integer
\(b\) prime to the residue characteristic of \(\lambda\)
that carries each of the adapted and standard lattices
into the other. If \(D_P\) bounds its total degree, then
\(b^{D_P}P\) has integral coefficients in the standard
letters. In degree \(a\) the lattice inclusions in the two directions give, on the tensors in the source,
\[
 b^a L_{\rm ad}^{\otimes a}\subseteq L_{\rm std}^{\otimes a},
 \qquad b^a L_{\rm std}^{\otimes a}\subseteq L_{\rm ad}^{\otimes a}.
\]
Dualizing reverses the inclusions and again costs at most \(b^a\) on coefficients. The same coefficient substitution in the
polynomial realizes these inclusions after imposing the
harmonic equations. For homogeneous degree \(a\le D_P\), a denominator cleared by \(b^a\) is also cleared by the same power for the whole polynomial, since
\[
  b^{D_P}=b^{D_P-a}b^a,\qquad D_P-a\ge0.
\]
 Homogeneity and the harmonic
equations are preserved. The prescribed highest-weight coefficient map gives a lower Fock vector for this
polynomial in the standard integral model.

The highest-weight coefficient map is defined over the common row field, identified with the
coefficient field of the lower vector through its chosen
embedding. Transport into the adapted coefficient lattice
can cost another power of \(b\). To see the support of
this denominator, at every place over the residue
characteristic the two source lattices coincide, so the
integral symmetric tensors and their differential
conditions coincide there as well. Outside those places
only finitely many primes require clearing. Descent by
the coefficient isomorphism and the integral highest-weight
map then supply a lower vector \(y\) for \(P\) and an
integer \(s_1\notin\lambda\) with
\[
                 s_1 y\in\mathcal C(h).
\]
Here \(\mathcal C(h)\) is the prescribed integral coefficient lattice, written in the lower Fock coordinates. The
product of the two denominator clears is allowed in
\(s_1\). This step uses the harmonic lattice on the
adapted source, and the coefficient lattice obtained from
the chosen highest-weight model.
\label{cor:paired-off-prime-trace-3}

\proofstep{The adelic theta evaluation.}
We describe the adelic evaluation used for this generator.
Let \(\chi_{\rm tuple}\) be the run's tuple character and
\(\chi_{\rm aux}\) its finite-order auxiliary character.
For the base kernel the splitting character is
\[
                    \chi_{\rm split}=\chi_{\rm tuple}\chi_{\rm aux}^{-1}.
\]
It is trivial on real-field ideles and has the prescribed values of the tame determinant character over the selected half of \(p_\circ\).
The type residues are the prescribed full-rank ones. The half
places are disjoint from their conjugates and from the
type places. The self-dual frames give the base kernel
as a product of a vacuum and the indicated residue
functions. Write \(\Phi_{\rm fr}\) for the finite vector
with free lower block at the twist places and the
specified full-rank residue matrices at the type places.
These full-rank residues are determined by the adapted
orbit's label. At the twist places the upper \(r\)-by-\(r\)
block is invertible and receives the inverse determinant
character.
\label{cor:paired-off-prime-trace-4}

For a finite vector \(\Phi\), let
\(S_\Phi(h,y;\mathbf w)\) be its single adelic theta series
with coefficient 
\[
  y_{\C}=c'_0(y)
\]
 and the tuple \(\mathbf w\) of
lowest target vectors. Here \(c'_0\) is the coefficient realization for this choice of levels. It sends the algebraic harmonic vector
to the complex coefficient representation using the fixed embedding.
Thus \(S_\Phi(h,y;\mathbf w)\) reads \(y_{\C}\) in the series,
and clearing \(s_1 y\) at \(h\) gives the coefficient
\(s_1 y_{\C}\) in the realized lattice. Denote the adelic realization of the coefficient row by \(\mathcal A(-)(\mathbf w)\). Let \(f\) be the single adapted short rule. With the
fixed choice of \(\Phi_{\rm fr}\) the extremal comparison
is the identity of adelic forms
\[
          \mathcal A(f)(\mathbf w)=c\,S_{\Phi_{\rm fr}}(h,y;\mathbf w),
          \qquad c\in K,\quad s_2c\in O_{K,\lambda},\quad s_2\notin\lambda .
\]
The coefficient and target tuple in this identity are the
ones selected above.
In this comparison one unfolds the geometric class sum
using the fixed adelic lift.

At the extremal tuple the
coefficient map differentiates exactly the letters of
the short polynomial, with the polarization sign and the
common frame and cover factors. The base-tuple class sums
span the cusp space, so the frame comparison applies to the sections defined by short theta rules on the whole row. To obtain that
span, choose a residue point using the surjective row
pairing at each split half-level place and use the
Chinese remainder theorem.
For the disjoint prime-power half-level conditions, choosing the local residue point amounts to using the surjection
\[
  {\cal O}_E\longrightarrow\prod_{\mathfrak p\ {\rm used}}{\cal O}_E/\mathfrak p^{n_{\mathfrak p}} .
\]
 The stable short span can
then be read at that point. Each raw term is moved to a
split-supported self-dual class by the short theta
transformation. The one-class kernel formula there gives
the geometric class sum times the frame factors. These
factors are nonzero, and linear independence of the cusp forms
descends the resulting span to the construction field.
\label{cor:paired-off-prime-trace-5}

\proofstep{Correct the level character in the adelic evaluation.}
The central level translation absorbs the difference
between the two frame normalizations. The base phase is
one, and the quotient of the two cover values and the
remaining central normalization are roots of unity in the field of values.
Thus if \(\xi_1,\xi_2\) denote the indicated normalization factors, their contribution at the chosen place is
\[
  v_\lambda(\xi_1\xi_2)=v_\lambda(\xi_1)+v_\lambda(\xi_2)=0 .
\]
 Their multiplication preserves the
local integrality test on \(c\). When the class
scalar is nonzero its expression, up to the preceding
units, is
\[
    B(\Lambda)^{-1}\Psi(\mathfrak a)
       [\Lambda:\Lambda\cap\Lambda_{\rm st}]^{r/2}.
\]
Here \(B(\Lambda)\) is the base genus scalar,
\(\Lambda_{\rm st}\) is the standard lattice, and
\(\Psi(\mathfrak a)\) is the unramified ideal value
of the combined tame character at the moved class.
Choose the determinant representative with
\(\mathrm N\mathfrak a\) prime to the residue
characteristic, using the local equality of the
adapted and standard lattices there.

In fact the determinant
lattice is the quotient of \(\mathfrak a\) by its conjugate,
and the two integral ideals are coprime. Order zero in
the determinant lattice at all places over the residue
characteristic forces order zero in both ideals there.
The local genus comparison gives \(B(\Lambda)\) order zero
at the same places, by comparing the theta normalization
in the common integral frame.
Consequently both directions of the genus-scalar comparison at such a place have order zero:
\[
  v_w\big(B(\Lambda)^{-1}\big)=-v_w\big(B(\Lambda)\big)=0 .
\]
 Its inverse becomes integral
after multiplication by an integer prime to that
characteristic. For the tame ideal value, take a power
of \(\mathfrak a\) that is generated by an element
congruent to one modulo the character conductor, using
the finite ray class group. Its character value is then
a product of integral powers of the embeddings of that
generator.

The generator divides a power of
\(\mathrm N\mathfrak a\) as an algebraic integer, so a
power of the norm clears the negative exponents.
Integrality descends from a positive power by
transitivity of integrality.
In particular, if \(a^n\) is integral, then \(a\) satisfies the monic equation
\[
  X^n-a^n=0.
\]
 The square of the positive
index factor is a rational integer. Their product is
therefore integral after multiplication by one integer
\(s_2\notin\lambda\). The computation takes place after extension to the chosen complex embedding, and integrality descends to \(K\) by
the injective field embedding. If the coefficient
is zero use \(s_2=1\).
The determinant-one condition on \(U\) makes the
remaining finite weight character trivial on the level
translates. Thus the same \(c\) applies to each term
\(f'_u\), with the unitary level action by \(u^{-1}\) on
the target series.
\label{cor:paired-off-prime-trace-6}

\proofstep{Extending the identity to all finite vectors.}
Both sides of this identity are multilinear in the
lowest target vectors and equivariant separately at
each archimedean place. The compact-group translates of
the nonzero extremal vector span the lowest type at
that place, by irreducibility of the lowest type.
Translate and extend linearly in each variable. Taking
the \(\psi\)-weighted sum gives
\begin{equation}\label{eq:off-prime-trace-identity}
 \mathcal A(t')(\mathbf w)
    =\sigma(c)\, T_{\psi}^{\mathrm{ad}}
                   S_{\Phi_{\rm fr}}(h,y;\mathbf w)
 \quad\text{for all }\mathbf w,
\end{equation}
where the adelic trace uses the same action and summation
as the row trace. The embedding \(\sigma\) here is the
one fixed for the adelic evaluation. The form coordinates
are determined by these tests. Indeed the smooth
archimedean module is generated by its lowest types, and
their translates on the target recover the full
equivariant form. The coordinate map sends its basis to linearly independent forms, hence is injective. In particular the vanishing test for the form coordinates reads
\[
  c=0\quad\Longrightarrow\quad t'=0.
\]
\label{cor:paired-off-prime-trace-7}

Finally turn the free-block vector into a translate
of the base kernel \(\Phi_{\rm b}\). At a chosen twist
place scale the lower \(n_V-r\) rows in the split model
by an element of valuation one, with the inverse adjoint
on the conjugate half. The upper \(r\) rows are fixed.
The original kernel imposes zero lower residues. After
the substitution this is automatic on integral entries.
Conversely, if the scaled matrix is integral with the
required zero lower residues, every entry before scaling
is integral.
For a lower-block entry \(a\) and local uniformizer \(\varpi\), the support equivalence is
\[
  \varpi a\in\varpi{\cal O}_v\quad\Longleftrightarrow\quad a\in{\cal O}_v .
\]
 The upper block determinant and its
character are preserved. This proves the equality of
the two support tests including the integrality test.

At a type prime, complete the prescribed \(r\) columns to an
invertible \(n_V\)-by-\(n_V\) matrix \(B\) over the residue field.
Since \(n_V>r\), column \(r+1\) is among those added in the completion.
Multiplying that column by \((\det B)^{-1}\) gives determinant one and
preserves all \(r\) prescribed columns. Lift the change to the split unitary
stabilizer by using the inverse adjoint at the conjugate
place. It fixes the vacuum and the coefficient lattice
there. The type and twist operations have disjoint
supports. They give the identity of finite vectors
\[
      \Phi_{\rm fr}=\sigma(c')\,\omega(k,1)\Phi_{\rm b}
\]
with a finite source translate \(k\) and nonzero \(c'\in
K\). In the split Weil formula, the determinant factors
of the twist operation use the local values of
\(\mu_0\chi\eta^{-1}\) at the chosen uniformizers.

\proofstep{The local determinant valuation.}
For each unit of determinant valuation in the forward
scaling, the norm factor is the inverse \(r\)th power
of the square root of the residue norm. The normalized
base character reading supplies the other inverse
square root. For residue cardinality \(Q\), the two scalar factors therefore combine as
\[
  Q^{-r/2}Q^{-1/2}=Q^{-(r+1)/2},\qquad (r+1)/2\in\mathbb Z.
\]
The translate uses the inverse scalar. The tame
character values and their inverses are integral
after multiplication by a power of \(p_\circ\), by
the same unramified ideal argument for the base
characters and finite order for the auxiliary
character.

At \(w\nmid p_\circ\) the two clearing conditions on a nonzero tame value \(\alpha\) imply
\[
  v_w(\alpha)\ge0,\qquad v_w(\alpha^{-1})\ge0\quad\Longrightarrow\quad v_w(\alpha)=0 .
\]
 The supported residue norms are powers
of \(p_\circ\). The type operation has determinant one.
The coefficient lattice also changes only at
\(p_\circ\), where a large enough power carries the
required coefficient lattices into one another. Hence
there is an \(m\ge0\) such that \(p_\circ^m c'\) is
integral and \(p_\circ^m s_1 y_{\C}\) is admissible at
\(hk\). Increasing \(m\) preserves both properties.
\label{cor:paired-off-prime-trace-8}

By the Weil representation law, the series with
\(\omega(k,1)\Phi_{\rm b}\) at \(h\) is the series of
\(\Phi_{\rm b}\) at \(hk\). For \(c\ne0\) let \(e\) be the
form coordinate of the latter kernel trace with
coefficient \(p_\circ^m s_1y_{\C}\). Equation
\eqref{eq:off-prime-trace-identity} and the determining
property of the tests identify it with
\[
               e=\frac{p_\circ^m s_1}{cc'}\,t'.
\]
It belongs to \(J\) by the admissibility just proved.

In the original orbit coordinates the identity reads
\[
       \bigl(p_\circ^{2m}s_1s_2\bigr)t
        =\bigl(\zeta^{-1}(s_2c)(p_\circ^m c')\bigr)e .
\]
The coefficient on the right is integral, and the
factor on the left is a unit at \(\lambda\). This proves
\(t\in J_\lambda\), also in the zero case. The argument
is uniform in the sense needed for the index comparison:
it incurs zero local order at every \(\lambda\) of
residue characteristic different from \(p_\circ\).
The clearing powers may depend on the generator and
the place. There is no such charge at the coefficient
place \(v\) because the run pin chooses \(p_\circ\) from
the other rational prime. Pull back along the common
form-coordinate isomorphism to get the inclusion on
the preprojected row. Applying the good projector
gives the inclusion on the good row.
\label{cor:paired-off-prime-trace-9}
\end{proof}
